\section{Full permanent multiplicities for a thin third row}

\subsection{O-021: a transverse permanent jet and padding theorem}
Let $r\geq3$, $n\geq m$, and let $\lambda\vdash nd$ satisfy
\[
       \ell(\lambda)\leq r,\qquad \lambda_3\leq2,\qquad
                    m\geq\max\{r,\lambda_2+1\}.
\]
Then, with the same internal padding as the supplied IP and
BIP papers,
\[
                    q^m_\lambda(d[n])
                         =K_n(\lambda)=a_\lambda(d[n]).
                                                               \tag{O-021.1}
\]
There is no upper bound on the outer degree $d$.
The statement determines all copies of these types on the
permanent side, rather than only their occurrence.

\paragraph{The unpadded permanent has the same transverse differential.}
For any $s\geq r-1$, use the variables
$x,y,z_1,\ldots,z_R$, $R=r-2$, and the difference-labeled
matrix of O-020, but change the signs on its upper-triangular
entries:
\[
        A_0=\operatorname{diag}(\ell_1,\ldots,\ell_s)+U,
        \qquad \ell_i=y-a_i x,\qquad
        U_{ij}=\zeta_i-\zeta_j\quad(i<j).
\]
Here $\zeta_i=z_i$ for $i\leq R$ and zero otherwise,
and the $a_i$ are distinct. The underlying directed graph
is acyclic. Consequently
$\operatorname{per}A_0=\Delta=\prod_i\ell_i$.

The permanent minor obtained by deleting row $j$ and
column $i$, with $i<j$, is the sum over directed paths
from $i$ to $j$ of the path's edge product multiplied
by the diagonal entries at its unused vertices.
Indeed every term of that minor consists of this path
and loops; an additional nontrivial cycle is impossible
in an acyclic graph. Conversely each path and its unused
loops gives exactly one minor term. In particular the
minor equals
\[
        (\zeta_i-\zeta_j)\frac{\Delta}{\ell_i\ell_j}
               +\text{terms of transverse degree at least two}.
                                                               \tag{O-021.2}
\]
The diagonal minors are $\Delta/\ell_i$.
This is an identity with weighted paths; negative coefficients
inside the difference labels cause no difficulty.

Vary diagonal entries by $x,y,z_a$, and lower-triangular
entries by linear forms in $z$. After deleting transverse
degree at least three, their permanent differentials are
exactly the selected determinant differentials in O-020.
The quadratic incidence isomorphism and the added-vertex
argument therefore apply unchanged. The permanent map
from all $s$-by-$s$ linear pencils dominates the full space
\[
 \operatorname{Sym}^s\mathbb C^2
 \ \oplus\ \bigl(\operatorname{Sym}^{s-1}\mathbb C^2
                                      \otimes\mathbb C^R\bigr)
 \ \oplus\ \bigl(\operatorname{Sym}^{s-2}\mathbb C^2
                               \otimes\operatorname{Sym}^2\mathbb C^R\bigr).
                                                               \tag{O-021.3}
\]
Its higher transverse coefficients are unrestricted.

\begin{proof}[Proof of the full multiplicity equality]
If $a_\lambda=0$, both quotient multiplicities vanish,
so assume it is positive. One-row types have
$q=K=a=1$, since the highest-weight function $F(e_1)^d$
is nonzero on both closures. Thus assume $\lambda_2>0$.
Then $\ell(\lambda)\leq d$.
Set $s=m-1$ and $b=\lambda_2$. The hypotheses give
$s\geq r-1$ and $s\geq b$. By F-014, the complete
ambient highest-weight multiplicity spaces at inner degrees
$s$ and $n$ have the same dimension:
\[
 a_\mu(d[s])=a_\lambda(d[n]),\qquad
             \mu=(sd-|\bar\lambda|,\bar\lambda).
                                                               \tag{O-021.4}
\]
Here $\bar\lambda=(\lambda_2,\lambda_3,\ldots)$.
The first part of $\mu$ is at least its second part,
because $|\bar\lambda|\leq(d-1)b$ and $s\geq b$.

The padding lift $T$ of O-012.b is injective from the
entire highest-weight space of weight $\mu$ into that
of weight $\lambda$. Equation (O-021.4) makes it an
isomorphism of these multiplicity spaces.
For the actual padding map
$P(F)=x^{n-s}F$, its pullback on $T$ is the diagonal
operator from the differential pairing in O-012.b:
\[
 P^*T=\operatorname{Sym}^d(\mathcal D),\qquad
 \mathcal D(e^\alpha)
        =\frac{(\alpha_1+n-s)!}{\alpha_1!}\,e^\alpha,
                 \qquad |\alpha|=s.                   \tag{O-021.5}
\]
Every diagonal entry is nonzero in characteristic zero.
The operator also preserves every monomial support condition.
It need not preserve the highest-weight subspace; that
additional property is not used.

By F-016 at its third-row cut, every polynomial in the
highest-weight space of weight $\mu$ uses only coefficient
generators of transverse degree at most two, since
$\mu_3=\lambda_3\leq2$. Thus the same is true after
applying the invertible operator in (O-021.5).
Any nonzero linear combination of the lifted copies
therefore pulls back to a nonzero polynomial in the
coordinates of (O-021.3). The dominant permanent map
proved above cannot annihilate that polynomial.
This proves independence of every lifted copy on the
padded permanent-pencil family.

Finally that entire family belongs to the required
internally padded permanent closure. For an arbitrary
$s$-by-$s$ pencil $A$, use the $m$-by-$m$ block matrix
\[
                         \widetilde A=\operatorname{diag}(x,A).
\]
Its distinguished entry is exactly $\widetilde A_{11}=x$,
and
\[
 \widetilde A_{11}^{\,n-m}\operatorname{per}\widetilde A
                      =x^{n-s}\operatorname{per}A.
                                                               \tag{O-021.6}
\]
These are linear endomorphism substitutions in
$X_{11}^{n-m}\operatorname{per}_m$. They belong to its
orbit closure because invertible endomorphisms are dense.
The variable space fits: $r\leq m\leq m^2\leq n^2$.
The independence just proved forces the entire permanent
ideal highest-weight space of type $\lambda$ to vanish.
Hence $q^m_\lambda=a_\lambda$.
O-020 also gives $K_n=a_\lambda$, since $n\geq m\geq r$.
This proves (O-021.1), including $n=m$.
\end{proof}

\paragraph{A strict four-copy example below every exact support-cover size.}
For a downward-closed support $\Gamma$ containing
$x_2^b$ and $x_i^2$ for every $3\leq i\leq r$, with
$b\geq2$, the two-tree invariant of F-018 satisfies
\[
                            \kappa(\Gamma)\geq b+r-2.
                                                               \tag{O-021.7}
\]
To prove this, let $P,Q$ be any admissible pair.
A decomposition of the pure power $b e_2$ forces
$a e_2\in P$, $a'e_2\in Q$ with $a+a'\geq b-1$.
Downward closure supplies both full axis chains, so these
alone contribute at least $a+a'+1\geq b$ to
$|P|+|Q|-1$. For $2e_i$, at least one of the two
endpoints has a positive $i$-th coordinate: the coefficient
edge can contribute at most one such unit. Downward closure
then forces $e_i\in P\cup Q$. These $r-2$ distinct
off-axis vertices are additional to the two chains.
This proves (O-021.7).

Take the concrete values
\[
       m=9,\quad n=10,\quad d=r=6,\qquad
                    \lambda=(44,8,2,2,2,2).
                                                               \tag{O-021.8}
\]
O-021.1 applies because $m=\max\{6,8+1\}$.
The exact near-rectangle formula F-013 gives
\[
      a_\lambda(6[10])=p_6(6,8)-p_5(6,8)=11-7=4.
\]
Thus
\[
                        q^9_\lambda(6[10])
                                =K_{10}(\lambda)=4.   \tag{O-021.9}
\]
Its nested support is
$\Gamma=\{\gamma:|\gamma|\leq8,\
                  \gamma_3+\cdots+\gamma_6\leq2\}$.
It contains the powers required by (O-021.7), so
$\kappa(\Gamma)\geq12>n>m$. Since the total-degree
cap is already eight, this is the same support at base
degrees $8,9,10$. Every exact two-tree support-cover
criterion of F-018 therefore misses this example, while
the transverse dominant projection and the all-copy lift
determine its actual multiplicities. This is a limitation
of that exact-cover architecture, not a lower bound for
arbitrary determinant or permanent representations.

\paragraph{Geometric and integral scope.}
The weighted path minors and the block substitution
(O-021.6) are integral polynomial identities. The
transverse differential uses the nonzero Vandermonde
factors and the selected incidence map whose exact
two-torsion is computed in O-020.8. The padding pullback
also uses invertible factorial scalars in characteristic
zero. No equality of integral Schur decompositions after
arbitrary base change is asserted.
The torus support is the comodule grading of SGA~3 I,
Propositions~4.7.2--4.7.3, pp.~34--36; F-014 and F-016
use the Cauchy and exterior functors of Macdonald I(4.3),
p.~63, Appendix~A(5.3)--(5.6), pp.~155--156, and
SGA~6 VI, Theorem~3.3, pp.~392--393.
The argument retains the injective coordinate pullback
through the dominant permanent projection and the actual
padding map. It supplies multiplicity information beyond
the virtual lambda class, with no saturation inference.
